\subsection{ANGLES}

Since the field \(\mathbf{R}\) is ordered, we may orient the real number plane \(\mathbf{R}^2\) by taking \(e_1 \wedge e_2\) as positive bivector ( \(e_1, e_2\) being the vectors of the canonical basis). In the oriented real number plane \(\mathbf{R}^2\) (identified with \(\mathbf{C}\) in what follows) we can then define the angle \((\widehat{\Delta_1, \Delta_2})\) of an arbitrary pair of rays ( \(\Delta_1, \Delta_2\) ) with origin \(o\) (*). The set \(\mathfrak{A}\) of all angles has the structure of an abelian group (written additively) defined by

\[
(\widehat {\Delta_ {1} , \Delta_ {3}}) = (\widehat {\Delta_ {1} , \Delta_ {2}}) + (\widehat {\Delta_ {2} , \Delta_ {3}}),
\]

so that, in particular, \((\widehat{\Delta_1, \Delta_1}) = 0\) and \((\widehat{\Delta_2, \Delta_1}) = -(\widehat{\Delta_1, \Delta_2})\) .

The flat angle \(\varpi\) is the solution \(\neq 0\) of the equation \(2\theta = 0\) in \(\mathfrak{A}\) ; it is the angle which the negative real semi-axis makes with the positive real semi-axis.

If z is any non-zero complex number, the amplitude (or argument) of z, denoted by \(\operatorname{Am}(z)\) , is the angle which the ray through z with origin o makes with the positive real semi-axis. The mapping \(z \to \operatorname{Am}(z)\) is a homomorphism of the multiplicative group \(C^{*}\) onto the additive group A, and therefore we have

\[
\operatorname{Am} \left(z z ^ {\prime}\right) = \operatorname{Am} (z) + \operatorname{Am} \left(z ^ {\prime}\right) \quad \text { and } \quad \operatorname{Am} (\bar {z}) = \operatorname{Am} \left(z ^ {- 1}\right) = - \operatorname{Am} (z).
\]

The angle \(\delta = \operatorname{Am}(i)\) is called the positive right angle; it is one of the solutions in A of the equation \(2\theta = \varpi\) , the other one being \(-\delta = \delta + \varpi\) .

The homomorphism \(z \to \operatorname{Am}(z)\) , restricted to the subgroup \(\mathbf{U}\) of \(\mathbf{C}^*\) , is an isomorphism of the group structure of \(\mathbf{U}\) onto that of \(\mathfrak{A}^{(**)}\) ; if we use this isomorphism to transport the topology of U to the group A, the latter group becomes a compact topological group, and the mapping \(z \to \operatorname{Am}(z)\) of \(C^{*}\) onto A is a strict morphism of the topological group \(C^{*}\) onto the topological group A.

Let us denote by \(\theta \to f(\theta)\) the isomorphism of \(\mathfrak{A}\) onto \(\mathbf{U}\) which is the inverse of the isomorphism \(z \to \operatorname{Am}(z)\) of \(\mathbf{U}\) onto \(\mathfrak{A}\) . By definition \(\mathcal{R}(f(\theta))\) is denoted by \(\cos \theta\) and is called the cosine of the angle \(\theta\) ; \(\mathfrak{J}(f(\theta))\) is denoted by \(\sin \theta\) and is called the sine of the angle \(\theta\) . These functions are continuous on the topological group \(\mathfrak{A}\) , and satisfy the following relations (loc. cit.), which are immediate consequences of the definitions above:

\[
\begin{array}{r l} \cos o = 1, & \sin o = 0, \quad \cos \varpi = - 1, \quad \sin \varpi = 0, \\ \cos (- \theta) = \cos \theta , & \sin (- \theta) = - \sin \theta , \\ \cos (\theta + \theta^ {\prime}) = \cos \theta \cos \theta^ {\prime} - \sin \theta \sin \theta^ {\prime}, \\ \sin (\theta + \theta^ {\prime}) = \sin \theta \cos \theta^ {\prime} + \sin \theta^ {\prime} \cos \theta , \\ \cos^ {2} \theta + \sin^ {2} \theta = 1. \end{array}
\]

By definition, the tangent of an angle \(\theta \in \mathfrak{A}\) is defined, whenever \(\cos \theta \neq 0\) , to be \(\sin \theta / \cos \theta\) (loc. cit.) and is denoted by \(\tan \theta\) ; it is a continuous function, which extends by continuity to \(\tilde{\mathbf{R}}\) (Chapter VI, § 3, no. 4) by taking the value \(\infty\) for the angles \(\delta\) and \(-\delta\) . We have \(\tan (\theta + \varpi) = \tan \theta\) . The cotangent of \(\theta\) , denoted by \(\cot \theta\) , is the element of \(\tilde{\mathbf{R}}\) equal to \(1 / \tan \theta\) .

Note that, if \(\operatorname{Am}(z) = \theta\) , we have \(z = |z| (\cos \theta + i \sin \theta)\) ; this expression is called the trigonometric form of the complex number \(z \neq 0\) .

